\paragraph*{Record, report and checks.} All values go into one record with their sources. A local language model
writes the report from that record alone, and two rule-based programs check every number and every direction against
it. Every score, measurement and statement sits in the record under a fixed key. The language model, Qwen3.8-27B
served with vLLM \citep{kwon2023} under greedy decoding, runs on the computer that runs PRISM, so the record never
leaves it. Decoding is constrained to a JSON schema in which every sentence carries its section and the keys of the
record fields it uses. The model sees no images, has no tools and does no arithmetic. The report page follows the
content of the BT-RADS follow-up template \citep{weinberg2018}, with the measurements in tables and a short narrative,
and without the template's category or advice, which PRISM does not give (Supplementary Section~S8). Because the
supplied volumes are noisy (Section~\ref{sec:3.3}), the report calls them indicative.

The checker, frozen before any evaluation report was generated, rejects a report that cannot be parsed, cites a field
that does not exist, states a number found in no cited field, omits a raised flag, misstates the conformal set or the
agreement, omits the card's tier where evidence is thin, leaves out a score, quotes a performance figure absent from
the card, or offers a diagnosis, a response category, advice or words of certainty. A rejected report is not
repaired. Such rule-based checks compare generated text with a source record \citep{zeba2026}, as extractive
evaluations of data-to-text systems do with a learned extractor \citep{wiseman2017}. Every LUMIERE report failed the
checker at the first evaluation, because a sentence about the rater's three-month flag stated a number held in no
field. We revised the record fields and the prompt once, on development patients, and generated the evaluation
reports again. Because the checker accepts a number that equals a number in any cited field, it accepts a sentence
that gives the FLAIR volume as the enhancing volume while citing both fields, or that writes ``larger'' where the
record says ``smaller''.

% Sources: METHODS_RESULTS.md M3; METHODS_PLAN.md M3; CHECKER_MUTATION.md.
The verifier closes these gaps with two rules that bind a claim to what it refers to. A sentence is split into clauses
at commas, semicolons, colons and the word ``and'', and each number and direction word in a clause is assigned to the
nearest tumour region or score named in it. The first rule requires every such number to equal, up to its written
rounding, a cited value of that region or score. The second requires every direction word on a fixed list (larger,
increased, grew, rose, higher, and their opposites) to agree with the cited direction of its region. As first
written, the rules read the region only from the names of cited fields and rejected every evaluation report, because
most fields are sentences such as ``enhancing tumour 3.5\,cc larger than at the previous scan''. We let the rules read
the region from the text of a cited field as well, tested this change only on synthetic records and the 48 MU and
LUMIERE reports used to develop the prompt, froze it, and ran it again on the evaluation reports. To measure the
verifier we planted one error at a time in correct reports, with 24 kinds of error, as error injection is used to
evaluate factuality checks \citep{gabriel2021,rao2025rexerr}. Three binding errors were held out of rule writing.
Because we wrote both the rules and these errors, a separate language model that had not seen the rules then wrote
further errors in its own words, with correct paraphrases as controls (exploratory). Both programs test whether a
report agrees with its record. Neither tests whether the record, or the report, is clinically correct.

